\chapter{Reaksi Wittig dan Reaksi Pembentuk C=C Lainnya}\label{ch:b2:wittig}

Lalat rumah betina menarik lalat jantan dengan hidrokarbon berkarbon 23 yang
membawa satu ikatan rangkap saja, di antara karbon 9 dan 10, dengan geometri
(\textit{Z}). Untuk membuatnya dalam jumlah banyak, seorang kimiawan harus
menempatkan satu ikatan rangkap itu tepat di sana, dan dengan geometri itu:
sebuah eliminasi akan menyebarkannya ke beberapa posisi dan menghasilkan
kedua geometri. Bab ini membangun ikatan rangkap dengan cara lain, yaitu
dengan menyambungkan dua potongan pada gugus karbonil, sehingga ikatan
rangkap itu muncul di tempat \ce{C=O} semula.
% ledger: use:muscalure

\begin{recall}
\Cref{ch:b2:enolates-aldol}: \omterm{def:b2:enolates-aldol:aldol}{kondensasi aldol}. \Cref{ch:b2:alkene-redox}:
hidrogenasi alkuna menjadi alkena (\textit{Z}) dengan katalis yang
diracuni. Jilid Tahun 1: E1, E2, dan aturan Zaitsev, substitusi
bimolekular, pereaksi Grignard dan pembalikan kepolaran, deskriptor
(\textit{Z})/(\textit{E}). Jilid sekolah: ekonomi atom.
\end{recall}

\section{Garam fosfonium dan ilida}

\begin{definition}[Ilida fosfonium]\label{def:b2:wittig:ylide}
Sebuah \emph{garam fosfonium}\index{garam fosfonium} $\mathrm{R_3P^+{-}CH_2R'\ X^-}$
dibuat melalui substitusi bimolekular haloalkana oleh fosfina, biasanya
trifenilfosfina \ce{Ph3P}. Melepaskan sebuah hidrogen dari karbon yang
terikat pada fosfor menghasilkan \emph{ilida fosfonium}\index{ilida fosfonium},
yaitu molekul netral dengan muatan-muatan berlawanan yang bersebelahan,
$\mathrm{Ph_3P^+{-}CH^-{-}R'}$, yang juga ditulis $\mathrm{Ph_3P{=}CHR'}$.
Sebuah \emph{ilida terstabilkan}\index{ilida terstabilkan} membawa pada
karbon itu gugus penarik elektron (ester, keton, \omterm{def:b2:amines:nitrile}{nitril}) yang
mendelokalisasikan muatan negatifnya.
\end{definition}

\begin{omfigure}
\setchemfig{atom sep=1.9em}
\schemestart
\chemfig{Ph_3P} \+ \chemfig{CH_3-I}
\arrow{->[S$_\mathrm{N}$2]}[,0.9]
\chemfig{Ph_3P^{+}-CH_3}\;\chemfig{I^{-}}
\arrow{->[\ce{BuLi}][$-$ \ce{BuH}, \ce{LiI}]}[,1.3]
\chemfig{Ph_3P^{+}-\chemabove{C}{\scriptstyle -}H_2}
\arrow{<->}[,0.8]
\chemfig{Ph_3P=CH_2}
\schemestop
\par\vspace{6pt}

{\small Metilidenatrifenilfosforana, ilida yang paling sederhana.
Trifenilfosfina menggeser iodida dari iodometana; butillitium melepaskan
sebuah proton dari gugus metil. Kedua struktur di sebelah kanan adalah dua
cara menuliskan molekul yang sama: struktur dengan muatan terpisah lebih
dekat dengan kenyataan.}
\end{omfigure}

\begin{proposition}[Membuat ilida]\label{prop:b2:wittig:ylide-acidity}
\omterm{def:b2:wittig:ylide}{Garam fosfonium} jauh lebih asam pada karbon di sebelah fosfor daripada
alkana. Garam alkil fosfonium sederhana tetap memerlukan basa yang sangat
kuat (butillitium, natrium hidrida, natrium amida) dalam kondisi kering dan
inert; garam yang \ce{CH2}-nya juga membawa gugus karbonil dideprotonasi
oleh basa lemah (natrium karbonat, natrium hidroksida), dan \omterm{def:b2:wittig:ylide}{ilida
terstabilkan} yang terbentuk dapat diisolasi dan disimpan.
\end{proposition}

\begin{proof}[Argumen]
Fosfor positif di sebelah karbanion menstabilkannya melalui muatannya dan
melalui kemudahan terpolarisasinya atom fosfor yang besar: ilida sederhana
adalah karbanion yang terstabilkan, tetapi kurang terstabilkan daripada
\omterm{def:b2:enolates-aldol:enolate}{enolat}, karena itu diperlukan basa kuat. Karbonil yang bersebelahan
menyebarkan muatan negatif ke oksigen, seperti pada \omterm{def:b2:enolates-aldol:enolate}{enolat}
(\cref{prop:b2:enolates-aldol:alpha-acidity}), selain pada fosfor: asam
konjugasinya menjadi cukup asam untuk basa lemah, dan ilidanya, yang jauh
kurang basa, tidak lagi bereaksi dengan air atau udara.
\end{proof}

\section{Reaksi Wittig}

\begin{definition}[Reaksi Wittig]\label{def:b2:wittig:wittig}
Yang disebut \emph{reaksi Wittig}\index{reaksi Wittig} antara \omterm{def:b2:wittig:ylide}{ilida
fosfonium} dan aldehida atau keton menghasilkan alkena dan trifenilfosfina
oksida:
\begin{equation*}
\mathrm{Ph_3P{=}CHR + R'CHO \longrightarrow R'CH{=}CHR + Ph_3P{=}O}.
\end{equation*}
Reaksi ini melalui sebuah \emph{oksafosfetana}\index{oksafosfetana}, yaitu
cincin beranggota empat yang terdiri atas fosfor, dua karbon, dan oksigen,
yang terbentuk ketika karbon ilida berikatan dengan karbon karbonil
sementara oksigen karbonil berikatan dengan fosfor.
\end{definition}

\begin{omfigure}
\setchemfig{atom sep=1.9em}
\schemestart
\chemfig{Ph_3P=CH-R} \+ \chemfig{O=CH-R'}
\arrow{->[adisi]}[,1.0]
\chemfig{Ph_3P?-[6]CH(-[4]R)-[0]CH(-[0]R')-[2]O?}
\schemestop

\medskip
\schemestart
\arrow{->[eliminasi]}[,1.1]
\chemfig{R-CH=CH-R'} \+ \chemfig{Ph_3P=O}
\schemestop
\par\vspace{6pt}

{\small \omterm{def:b2:wittig:wittig}{Reaksi Wittig}. Karbon ilida dan karbon karbonil saling berikatan
sementara oksigen berikatan dengan fosfor, dalam satu tahap: \omterm{def:b2:wittig:wittig}{oksafosfetana}.
Cincinnya lalu terbuka ke arah lain, dengan memutus ikatan \ce{P-C} dan
\ce{C-O}: ikatan \ce{C=C} terbentuk di antara kedua bekas karbon yang
bereaksi, dan ikatan
\ce{P=O} trifenilfosfina oksida terbentuk pada saat yang sama.}
\end{omfigure}

\begin{proposition}[Ikatan rangkap tepat di tempatnya]\label{prop:b2:wittig:regiospecific}
Dalam \omterm{def:b2:wittig:wittig}{reaksi Wittig}, ikatan \ce{C=C} yang baru menghubungkan bekas karbon
karbonil dengan bekas karbon ilida, dan tidak di tempat lain.
\end{proposition}

\begin{proof}[Argumen]
Satu-satunya ikatan yang terbentuk dan terputus adalah ikatan-ikatan cincin
beranggota empat: karbon ilida dengan karbon karbonil, lalu ikatan \ce{C-P}
dan \ce{C-O}. Tidak ada hidrogen yang berpindah dan tidak ada karbokation
atau karbanion yang terbentuk di sepanjang rantai, sehingga ikatan
rangkapnya tidak dapat bergeser. Sebaliknya, eliminasi memilih di antara
hidrogen-hidrogen pada beberapa karbon tetangga (aturan Zaitsev), dan reaksi
E1 dapat mengalami penataan ulang.
\end{proof}

\begin{proposition}[Gaya pendorong]\label{prop:b2:wittig:driving}
Pembentukan ikatan \ce{P=O} trifenilfosfina oksida membuat \omterm{def:b2:wittig:wittig}{reaksi Wittig}
sangat menguntungkan dan ireversibel.
\end{proposition}

\begin{proof}[Argumen]
Reaksi ini menukar satu ikatan \ce{C=O} dan satu ikatan \ce{P-C} (dalam
ilida, \ce{P=C} sebagian besar merupakan ikatan tunggal
\ce{P+-C-}) dengan satu ikatan \ce{C=C} dan satu ikatan \ce{P=O}. Ikatan
\ce{P=O} sebuah fosfina oksida sangat kuat, jauh lebih kuat daripada ikatan
\ce{P-C} yang hilang, dan \ce{C=O} yang hilang sebagian diimbangi oleh
\ce{C=C} yang terbentuk: neracanya jelas menurun. Fosfor adalah penampung
oksigen dalam reaksi ini.
\end{proof}

\section{Stereoselektivitas}

\begin{proposition}[Geometri alkena]\label{prop:b2:wittig:stereo}
Dengan aldehida, ilida tak terstabilkan (\ce{Ph3P=CHR}, R gugus alkil)
terutama menghasilkan alkena (\textit{Z}),
dalam kondisi bebas garam. \omterm{def:b2:wittig:ylide}{Ilida terstabilkan} (\ce{Ph3P=CH-COOR}, \ce{Ph3P=CH-COR})
terutama menghasilkan alkena
(\textit{E}). Ilida yang hanya distabilkan oleh gugus aril menghasilkan campuran.
\end{proposition}

\admitted

Ini adalah hasil pengamatan. Penjelasan yang lazim membaca \omterm{def:b2:wittig:wittig}{reaksi Wittig}
sebagai persaingan antara laju dan kestabilan. Dengan ilida tak
terstabilkan yang reaktif, \omterm{def:b2:wittig:wittig}{oksafosfetana} terbentuk dengan cepat dan secara
ireversibel, melalui keadaan transisi yang kurang sesak, yang kedua
substituennya berakhir pada sisi cincin yang sama; \omterm{def:b2:wittig:wittig}{oksafosfetana} itu lalu
runtuh dengan mempertahankan susunan tersebut: kendali kinetik menghasilkan
alkena (\textit{Z}). Dengan \omterm{def:b2:wittig:ylide}{ilida terstabilkan}, tahap pertama lebih lambat
dan reversibel; \omterm{def:b2:wittig:wittig}{oksafosfetana-oksafosfetana} saling berubah melalui bahan
awal, dan \omterm{def:b2:wittig:wittig}{oksafosfetana} dengan substituen-substituen pada sisi yang
berseberangan, yang kurang tegang, yang menang: alkena (\textit{E}), di
bawah kendali termodinamik.

\begin{omfigure}
\setchemfig{atom sep=1.9em}
\schemestart
\chemfig{Ph_3P=CH-CH_2CH_3} \+ \chemfig{O=CH-Ph}
\arrow{->}[,0.8]
\chemfig{Ph-[:-30]=[:30]-[:90]CH_2CH_3}
\schemestop

\medskip
\schemestart
\chemfig{Ph_3P=CH-COOEt} \+ \chemfig{O=CH-Ph}
\arrow{->}[,0.8]
\chemfig{Ph-[:-30]=[:30]-[:-30]COOEt}
\schemestop
\par\vspace{6pt}

{\small Dua ilida, satu aldehida. Atas: ilida propilidena yang tak
terstabilkan terutama menghasilkan (\textit{Z})-1-fenilbut-1-ena. Bawah:
\omterm{def:b2:wittig:ylide}{ilida terstabilkan} terutama menghasilkan etil
(\textit{E})-3-fenilpropenoat (etil sinamat). Trifenilfosfina oksida, yang
terbentuk dalam keduanya, tidak ditampilkan.}
\end{omfigure}

\section{Reaksi Horner--Wadsworth--Emmons}

\begin{definition}[Fosfonat dan reaksi HWE]\label{def:b2:wittig:hwe}
Yang disebut \emph{fosfonat}\index{fosfonat} adalah ester asam fosfonat,
$\mathrm{(RO)_2P({=}O){-}R'}$, dengan sebuah
ikatan \ce{P-C}. Fosfonat dibuat melalui reaksi Michaelis--Arbuzov antara
trialkil fosfit dan haloalkana,
$\mathrm{P(OEt)_3 + R'Br \to (EtO)_2P(O)R' + EtBr}$. Dalam
\emph{reaksi Horner--Wadsworth--Emmons}\index{reaksi Horner--Wadsworth--Emmons}
(HWE), anion fosfonat yang membawa gugus penarik elektron pada \ce{CH2}-nya
beradisi pada aldehida dan menghasilkan alkena serta anion dialkil fosfat.
\end{definition}

\begin{omfigure}
\setchemfig{atom sep=1.8em}
\schemestart
\chemfig{EtO-P(=[2]O)(-[6]OEt)-CH_2-COOEt}
\arrow{->[\ce{NaH}][$-$ \ce{H2}]}[,0.9]
\chemfig{EtO-P(=[2]O)(-[6]OEt)-\chemabove{C}{\scriptstyle -}H-COOEt}
\schemestop

\medskip
\schemestart
\arrow{->[\ce{PhCHO}]}[,0.9]
\chemfig{*6(=-=(-[:30]=[:-30]-[:30](=[:90]O)-[:-30]O-[:30]Et)-=-)}
\+
\chemfig{EtO-P(=[2]O)(-[6]O^{-})-OEt}
\schemestop
\par\vspace{6pt}

{\small \omterm{def:b2:wittig:hwe}{Reaksi Horner--Wadsworth--Emmons} antara trietil fosfonoasetat dan
benzaldehida. Natrium hidrida melepaskan hidrogen di antara \omterm{def:b2:wittig:hwe}{fosfonat} dan
ester; anionnya beradisi pada aldehida, dan zat antara beranggota empat
runtuh seperti dalam \omterm{def:b2:wittig:wittig}{reaksi Wittig}. Produknya adalah etil
(\textit{E})-sinamat; produk sampingnya adalah anion dietil fosfat, yang
larut dalam air.}
\end{omfigure}

\begin{proposition}[Mengapa kimiawan menyukai reaksi HWE]\label{prop:b2:wittig:hwe}
Reaksi HWE terutama menghasilkan alkena (\textit{E}), dan produk samping
fosfornya, garam yang larut dalam air, disingkirkan dengan pencucian air
sederhana; anion \omterm{def:b2:wittig:hwe}{fosfonat} juga lebih nukleofilik daripada \omterm{def:b2:wittig:ylide}{ilida
terstabilkan} yang bersesuaian dan juga bereaksi dengan keton.
\end{proposition}

\begin{proof}[Argumen]
Anion \omterm{def:b2:wittig:hwe}{fosfonat} adalah karbanion yang terstabilkan, sehingga adisinya
reversibel dan reaksinya berlangsung di bawah kendali termodinamik:
(\textit{E}), seperti pada \omterm{def:b2:wittig:ylide}{ilida terstabilkan}. Tidak seperti \omterm{def:b2:wittig:ylide}{ilida
fosfonium}, anion itu membawa muatan negatif penuh, yang tidak diimbangi oleh
kation di sebelahnya: anion itu lebih reaktif. Produk sampingnya,
\ce{(EtO)2PO2-}, adalah garam ionik, sedangkan trifenilfosfina oksida adalah
padatan organik netral yang harus dipisahkan dari produk dengan
kromatografi atau kristalisasi.
\end{proof}

\section{Memilih metode}

\begin{method}[Pemutusan Wittig]\label{met:b2:wittig:wittig-disconnection}
\begin{enumerate}
  \item Potonglah ikatan \ce{C=C} target; satu karbon menjadi karbon
        karbonil, yang lain menjadi karbon ilida (sebuah \omterm{def:b2:wittig:ylide}{garam fosfonium},
        dari sebuah halida).
  \item Dari kedua cara melakukannya, utamakan cara yang halidanya primer
        (substitusi bimolekular oleh
        \ce{Ph3P}) dan senyawa karbonilnya aldehida.
  \item Pilihlah pereaksinya menurut geometri yang diinginkan: ilida tak
        terstabilkan untuk (\textit{Z}), \omterm{def:b2:wittig:ylide}{ilida terstabilkan} atau \omterm{def:b2:wittig:hwe}{fosfonat}
        (HWE) untuk (\textit{E}).
  \item Periksalah bahwa tidak ada bagian lain molekul yang bereaksi dengan
        basa yang dipakai.
\end{enumerate}
\end{method}

\begin{method}[Memilih reaksi pembentuk C=C]\label{met:b2:wittig:choose-cc}
\begin{center}
\small
\begin{tabular}{@{}lll@{}}
\hline
Metode & Posisi \ce{C=C} & Geometri \\
\hline
Eliminasi E2 atau E1 & aturan Zaitsev, campuran & terutama (\textit{E}), campuran \\
\omterm{def:b2:enolates-aldol:aldol}{Kondensasi aldol} & terkonjugasi dengan \ce{C=O} & (\textit{E}) \\
Wittig, ilida tak terstabilkan & pada bekas \ce{C=O} & (\textit{Z}) \\
Wittig, \omterm{def:b2:wittig:ylide}{ilida terstabilkan}; HWE & pada bekas \ce{C=O} & (\textit{E}) \\
Alkuna, lalu katalis teracuni & pada bekas ikatan rangkap tiga & (\textit{Z}) \\
\hline
\end{tabular}
\end{center}
\medskip
Untuk menyambungkan dua potongan yang ukurannya mirip pada sebuah \ce{C=C},
\omterm{def:b2:wittig:wittig}{reaksi Wittig} dan kerabat-kerabatnya biasanya menjadi pilihan pertama;
metatesis olefin, yang menukar ujung-ujung dua alkena, dibahas dalam jilid
Tahun 3.
\end{method}

\begin{history}[Ilida yang menjadi pereaksi]
Georg Wittig, yang mempelajari ilida fosfor pada awal 1950-an, menemukan
bahwa metilidenatrifenilfosforana mengubah benzofenon menjadi
1,1-difeniletena dan trifenilfosfina oksida; bersama mahasiswanya Ulrich
Schöllkopf ia mengubah pengamatan itu menjadi metode umum pada tahun 1954.
Industri segera mengadopsinya, antara lain untuk vitamin A, dan Wittig
berbagi Hadiah Nobel Kimia 1979.
\end{history}

\begin{safety}{\ghs[5mm]{GHS02}\,\ghs[5mm]{GHS05}\,\ghs[5mm]{GHS07}\,\ghs[5mm]{GHS08}}
Larutan butillitium mudah terbakar dan korosif serta bereaksi hebat dengan
air; natrium hidrida melepaskan hidrogen dengan air. Trifenilfosfina
berbahaya dan korosif terhadap mata, trietil fosfit mudah terbakar dan
berbahaya; trifenilfosfina oksida berbahaya.
% ledger: ghs:BuLi, ghs:PPh3, ghs:triethyl-phosphite, ghs:Ph3PO
\end{safety}

\section{Latihan}

\begin{exercise}[$\star$]\label{exo:b2:wittig:1}
Tuliskan pembuatan ilida \ce{Ph3P=CHCH3} dari trifenilfosfina dan sebuah haloalkana.
\end{exercise}

\begin{exercise}[$\star$]\label{exo:b2:wittig:2}
Tuliskan produk-produk \omterm{def:b2:wittig:wittig}{reaksi Wittig} benzaldehida dengan ilida yang dibuat dari iodometana.
\end{exercise}

\begin{exercise}[$\star$]\label{exo:b2:wittig:3}
Golongkan sebagai terstabilkan, tak terstabilkan, atau terstabilkan aril:
\ce{Ph3P=CHCOOEt}, \ce{Ph3P=CHCH3},
\ce{Ph3P=CHCOCH3}, \ce{Ph3P=CHC6H5}. Manakah yang memerlukan butillitium?
\end{exercise}

\begin{exercise}[$\star$]\label{exo:b2:wittig:4}
Tuliskan produk trietil fosfonoasetat, natrium hidrida, dan propanal, beserta geometrinya.
\end{exercise}

\begin{exercise}[$\star\star$]\label{exo:b2:wittig:5}
Metilenasikloheksana diinginkan. Bandingkan jalur Wittig dari sikloheksanon
dengan dehidrasi terkatalisis asam 1-metilsikloheksan-1-ol.
\end{exercise}

\begin{exercise}[$\star\star$]\label{exo:b2:wittig:6}
Putuskan (\textit{E})-1,2-difeniletena (stilbena) dengan \omterm{def:b2:wittig:wittig}{reaksi Wittig}.
Mengapa jalur Wittig merupakan pilihan yang buruk untuk (\textit{E})-stilbena
murni, dan apa yang sebaiknya dipakai sebagai gantinya?
\end{exercise}

\begin{exercise}[$\star\star$]\label{exo:b2:wittig:7}
Propanal bereaksi dengan \ce{Ph3P=CHCH3} di satu pihak dan dengan
\ce{Ph3P=CHCOOEt} di pihak lain. Tuliskan produk-produknya dan geometri
utamanya.
\end{exercise}

\begin{exercise}[$\star\star$]\label{exo:b2:wittig:8}
Hitunglah ekonomi atom metilenasi sikloheksanon oleh \ce{Ph3P=CH2}.
\end{exercise}

\begin{exercise}[$\star\star$]\label{exo:b2:wittig:9}
Tuliskan reaksi Michaelis--Arbuzov trietil fosfit dengan etil bromoetanoat,
dalam dua tahap: substitusi oleh fosfor, lalu dealkilasi oleh bromida.
\end{exercise}

\begin{exercise}[$\star\star\star$]\label{exo:b2:wittig:10}
Usulkan sintesis 1,6-difenilheksa-1,3,5-triena dari (\textit{E})-but-2-enadial
\ce{OHC-CH=CH-CHO} melalui dua \omterm{def:b2:wittig:wittig}{reaksi Wittig}, dan sebutkan ilida mana dan
berapa banyak.
\end{exercise}

\begin{exercise}[$\star\star\star$]\label{exo:b2:wittig:11}
(\textit{E})-4-Fenilbut-3-en-2-on dapat dibuat melalui \omterm{def:b2:enolates-aldol:aldol}{kondensasi aldol} atau
melalui \omterm{def:b2:wittig:wittig}{reaksi Wittig}. Tuliskan kedua jalur itu dan bandingkan.
\end{exercise}

\begin{exercise}[$\star\star\star$]\label{exo:b2:wittig:12}
Usulkan dua jalur menuju (\textit{Z})-okt-4-ena, satu melalui \omterm{def:b2:wittig:wittig}{reaksi Wittig}
dan satu melalui alkuna, dan sebutkan produk samping yang ditinggalkan
masing-masing.
\end{exercise}

\section{Soal: Feromon melalui Reaksi Wittig}

\begin{problem}[{Soal akhir pekan --- retrosintesis sebuah alkena (Z), ilida
tak terstabilkan dan selektivitasnya, biaya berupa trifenilfosfina oksida,
dan alternatif melalui alkuna}]
\label{pb:b2:wittig:1}
Target: (\textit{Z})-trikos-9-ena, \ce{CH3(CH2)7CH=CH(CH2)12CH3}, feromon lalat
rumah pada pembuka bab. Massa molar (\unit{g/mol}): C 12.011, H 1.008,
O 15.999, P 30.974, Br 79.904.
% ledger: use:muscalure, aw:C, aw:H, aw:O, aw:P, aw:Br, ghs:nonanal, ghs:BuLi

\medskip
\noindent\textbf{Bagian I --- Retrosintesis.}
\begin{enumerate}
  \item Tuliskan rumus molekul target dan periksalah bahwa target itu alkena.
  \item Ikatan manakah yang dipotong oleh pemutusan Wittig?
  \item Tuliskan kedua pasangan yang mungkin (senyawa karbonil, \omterm{def:b2:wittig:ylide}{garam fosfonium}).
  \item Untuk pasangan nonanal + ilida, haloalkana manakah yang membentuk \omterm{def:b2:wittig:ylide}{garam fosfoniumnya}?
  \item Mengapa haloalkana primer cocok untuk tahap ini?
  \item Tuliskan pembentukan \omterm{def:b2:wittig:ylide}{garam fosfonium} itu.
\end{enumerate}

\medskip
\noindent\textbf{Bagian II --- Ilida dan geometri.}
\begin{enumerate}[resume]
  \item Basa manakah yang mendeprotonasi garam ini? Mengapa natrium
        hidroksida tidak memadai?
  \item Apakah ilida itu terstabilkan? Geometri apa yang diramalkan oleh hal itu?
  \item Tuliskan \omterm{def:b2:wittig:wittig}{oksafosfetana} yang terbentuk dengan nonanal.
  \item Mengapa reaksi harus dijalankan dalam keadaan kering dan di bawah gas inert?
  \item Mengapa ikatan rangkapnya berakhir tepat di antara karbon 9 dan 10?
  \item Apa yang akan dihasilkan oleh dehidrasi terkatalisis asam
        trikosan-9-ol sebagai gantinya?
\end{enumerate}

\medskip
\noindent\textbf{Bagian III --- Trifenilfosfina oksida.}
\begin{enumerate}[resume]
  \item Tuliskan persamaan keseluruhan dari ilida dan nonanal.
  \item Hitunglah massa molar target, nonanal, \omterm{def:b2:wittig:ylide}{garam fosfonium}, dan
        trifenilfosfina oksida.
  \item Hitunglah ekonomi atom tahap Wittig (ilida + aldehida).
  \item Bagaimana trifenilfosfina oksida dipisahkan dari sebuah hidrokarbon?
  \item Mengapa pembentukan trifenilfosfina oksida merupakan gaya pendorongnya?
  \item Mengapa reaksi HWE tidak cocok untuk target ini?
\end{enumerate}

\medskip
\noindent\textbf{Bagian IV --- Jalur alkuna.}
\begin{enumerate}[resume]
  \item Alkuna manakah yang, jika dihidrogenasi, menghasilkan target, dan
        dengan katalis apa?
  \item Bangunlah alkuna itu dari dek-1-una dan sebuah haloalkana: halida
        manakah, basa manakah?
  \item Mengapa hidrogenasi harus berhenti pada alkena, dan bagaimana
        katalisnya menjamin hal itu?
  \item Bandingkan kedua jalur: jumlah tahap, kendali geometri, produk samping.
  \item Berapakah ekonomi atom tahap hidrogenasi?
  \item Nyatakan massa trifenilfosfina oksida yang ikut dihasilkan per gram
        feromon melalui jalur Wittig.
\end{enumerate}
\end{problem}
